\documentclass[12pt]{article}

\usepackage{amsmath}
\usepackage{german}
\usepackage{graphicx}
\DeclareMathOperator\sgn{sgn}

%\graphicspath{{/home/ritz/Teaching/fcn/fig/}}

%----------------- layout ----------------------------------------

\oddsidemargin=0in
\textwidth=6in
\topmargin=-0.7in
\textheight=9.5in
\baselineskip=16pt
\font\mittel=cmr17
\font\roman=cmr12
\font\klein=cmr9
\font\mgross=cmr10
\newfont{\CAP}{cmcsc10 scaled 900}
\pagestyle{empty}
\parindent=0mm

%----------------- macros ----------------------------------------
\newcommand{\conc}[1]{[{\rm C}]_{\text{#1}}}
\newcommand{\units}[1]{\text{#1}}
\newcommand{\K}{\text{K}}
\newcommand{\Na}{\text{Na}}
\newcommand{\Cl}{\text{Cl}}
\newcommand{\Leak}{\text{leak}}
\newcommand{\water}{\text{H}_2\text{O}}
\newcommand{\tderiv}[1]{\frac{\mathrm{d}#1}{\mathrm{d}t}}
%-----------------------------------------------------------------


\begin{document}

\parbox{2cm}{
\includegraphics[width=1.8cm]{../bccnlogo.pdf}
}
\parbox{11cm}{
\begin{center}
\large HUMBOLDT-UNIVERSIT"AT \hskip 0.1 cm ZU \hskip 0.1 cm BERLIN
\vskip 0.1 true cm
\mgross BERNSTEIN CENTER FOR\\ COMPUTATIONAL NEUROSCIENCE
\end{center}
}
\parbox{2cm}
{
\hfill
\includegraphics[width=1.8cm]{../hublogo.pdf}
}

\vskip 0.6 true cm

\leftline{\CAP Humboldt-Universit"at zu Berlin 
\hfill Phone: 030/2093-9110}
\leftline{\CAP Philippstr. 13 House 6
             \hfill Fax: 030/2093-6771}
\leftline{\CAP \hfill webpage: http://www.bccn-berlin.de/}


\vskip 0.8 true cm
\centerline{\bf Models of Neural Systems, WS 2010/11}
\centerline{\bf Project 3: Adaptive exponential integrate-and-fire
model} 
\centerline{Project presentation and report submission: February,
15th, 2011}

\vskip 0.8 true cm

\noindent{\bf Background}

Detailed neuron models offer much biological realism and good
agreement with experimental findings. On the other hand, they are often
very complex, which makes parameter fitting, simulation and theoretical
analysis of such models very difficult. Therefore, simplified
neuron models are used to simulate neuronal dynamics. One class of such
models are based on the integrate-and-fire approach. How good are these
models in comparison with detailed conductance-based neuron models?
Here you will study the explanatory power of an adaptive exponential
integrate-and-fire neuron (aEIF).

\vskip 0.8 true cm

{\bf Problems}

\begin{enumerate}
    \item \textit{Literature review.} Study the paper of Brette and
        Gerstner,  J Neurophysiol 2005. Focus on the figures and  the method
        section. What is the hypothesis? What are the main results
        and conclusions? What are the main assumptions of the model?
        What are the strengths and weaknesses of the paper?

    \item Model the responses of a Hodgkin-Huxley model (HHM) to fluctuating
        synaptic inputs. To this end, add synaptic
        inhibitory and excitatory currents to the standard HHM:
        \begin{equation}
            I_\mathrm{syn}=g_\mathrm{syn}(E_\mathrm{syn}-V)
            \label{eq:syn_current}
        \end{equation}
        where $\mathrm{syn} \in  \left\{\mathrm{exc},
        \mathrm{inh}\right\}$ denotes excitatory or inhibitory
        currents. The respective reversal potentials are
        $E_\mathrm{exc}=0\text{ mV}$ and $E_\mathrm{inh}=-75\text{
        mV}$. The synaptic conductance $g_\mathrm{syn}$ is
        time-varying --  in first approximation, assume that
        they are independent Gaussian random numbers with mean
        $\overline{g}_\mathrm{syn}$ and standard deviation
        $\sigma_\mathrm{syn}$. Choose both parameters of the
        Gaussian distribution such that the neuron fires spontaneously
        with low firing rate (about 10 Hz). Consider both
        low-conductance (ratio of total mean conductance to leak conductance
        2:1) and high-conductance states (ratio 5:1).

    \item Implement the exponential integrate-and-fire neuron with
        adaption (aEIF) as described in the method section of
        Brette and Gerstner (2005).

    \item Fit the parameters of the aEIF model to the data generated from
        HHM according to the procedure described in Brette and
        Gerstner (2005).

    \item Compare the spike trains generated by the aEIF and HHM models
        to the same time-varying synaptic conductances in terms of
        missing and extra spikes (see Method section of Brette and
        Gerstner, 2005).

    \item Try to answer the following questions: Are simplified
        models a good description of the experimental data? How well
        can an aEIF  model reproduce the spike  trains generated by the HHM?
        How could the performance be improved? How well would you expect
        aEIF to describe real experimental data? Why?  A recent
        opinion on the issues can be also found in Gerstner and Naud,
        Science 2009.

\end{enumerate}

\vskip 0.8 true cm
{\bf Literature}
\vskip 0.4 true cm

R.~Brette and W.~Gerstner (2005). Adaptive integrate-and-fire model as an
effective description of neuronal activity.
\emph{J. Neurophysiol.} 94:3637-3642.\newline

W.~Gerstner and R.~Naud (2009). How good are neuron models?
\emph{Science} 326:379-380.\newline

\vfill
\centerline{\CAP Contact}
\CAP

\begin{tabular}{lll}
\underline{Urs Bergmann} & Phone: 2093-8924 & Email:
urs.bergmann(at)biologie.hu-berlin.de \\
Matthias Guggenmos & & Email: matthias.guggenmos(at)bccn-berlin.de \\
Richard Kempter \hfill & Phone: 2093-8925 \hfill & Email:
r.kempter(at)biologie.hu-berlin.de \\
Bartosz Telenczuk & Phone: 2093-8838 & Email:
b.telenczuk(at)biologie.hu-berlin.de \\
\end{tabular}vfill


\end{document}
